\documentclass[11pt,a4paper]{article}
\usepackage[T1]{fontenc}
\usepackage[utf8]{inputenc}
\usepackage{lmodern}
\usepackage{amsmath,amssymb}
\usepackage[margin=25mm]{geometry}
\usepackage[hidelinks]{hyperref}
\hypersetup{pdftitle={Gaussian product blocking: the tree-level coupling never moves},pdfauthor={navstokgap research working notes}}
\newcommand{\tightlist}{\setlength{\itemsep}{0pt}\setlength{\parskip}{0pt}}
\setlength{\emergencystretch}{3em}
\setcounter{secnumdepth}{0}
\begin{document}
\hypertarget{gaussian-product-blocking-the-tree-level-coupling-never-moves}{%
\section{Gaussian product blocking: the tree-level coupling never
moves}\label{gaussian-product-blocking-the-tree-level-coupling-never-moves}}

\textbf{Result, 2026-09-28 (Claude; written derivation, to be
refereed).} For the free (Gaussian) lattice gauge field in \(D\)
dimensions with heat time \(t\), one isotropic product-blocking step
\(a\mapsto2a\) maps the law of the fluxes to a Gaussian law whose
covariance at small coarse momentum \(K\) is

\[C'(K)=2^{4-D}\,t\,\bigl[P_{\rm cont}(\hat K)+O(K)\bigr]\]

on exact 2-forms, with \(O(K^2)\) in plaquette-centred phase conventions
(Theorem G1): the zero-momentum coefficient is exactly the heat time
\(\lambda_D(2a)^{4-D}\) of the coarse lattice. The statement holds for
any input covariance supported on exact 2-forms with the same
small-momentum form, so it iterates: along any number of steps, and for
the four directional halvings of one step in any order (Fubini), the
tree-level coupling of the Gaussian part is preserved exactly. In
\(1+3\) it does not run at all.

\textbf{Consequence for cell 4} (the
\href{four-dimensional-parallel-log.md}{four-dimensional note}, (15) and
the Remark after its telescoping paragraph). The endpoint term
\(c_{\rm in}-c_{\rm out}\) receives nothing from the zero-momentum part
of the Gaussian kernel. It arises only from the dependence of the
one-loop determinant on the kernel's shape (its \(O(K^2)\) and higher
terms) and from the generated cubic and quartic vertices. Along an
iteration the shape may drift; the average rate \(-2b_0\log2\) per step
needs only \(c(\zeta_n,\mathcal K_n)=o(n)\), so even a logarithmic drift
of shape parameters would be harmless.

\textbf{Caution.} Convergence of the full kernel under iterated sharp
product blocking is not claimed. Product blocking makes coarse fluxes
the fluxes of the fine field through sharp surfaces, whose variances in
the continuum carry boundary (perimeter-type) divergences; the
perfect-action literature smears the blocking to obtain a local fixed
point (\href{https://doi.org/10.1103/PhysRevB.11.3431}{Bell and Wilson
1975}; for gauge fields
\href{https://doi.org/10.1016/0550-3213(95)00678-8}{Bietenholz and Wiese
1996}; metadata).

\hypertarget{setting}{%
\subsection{1. Setting}\label{setting}}

Take \(G=\mathbb R\) (colour components decouple in the Gaussian theory)
on a periodic lattice of spacing \(a\), link variables \(A\), fluxes
\(\phi=dA\) and weight \(\prod_pe^{-\phi_p^2/(2t)}\). Modulo gauge,
\(\phi\) ranges over the exact 2-forms, with the flat Gaussian density
there, so its covariance is \(tP\), where \(P(k)\) is the orthogonal
projector onto exact 2-forms at lattice momentum \(k\) (\(k\ne0\); the
harmonic sector at \(k=0\) is set aside). Product blocking sets the
coarse link to the sum of the two fine links along it; by Stokes the
coarse flux through a coarse plaquette in the plane \((\mu\nu)\) is the
sum of the four fine fluxes inside it,
\((B\phi)_{\mu\nu}(X)=\sum_{a,b\in\{0,1\}}\phi_{\mu\nu}(2X+ae_\mu+be_\nu)\).
This is the Gaussian case of the pushforward of Proposition 1 of the
\href{series-parallel-gauge-refinement.md}{series/parallel note},
composed over the \(D\) directions.

\hypertarget{theorem-g1-and-its-proof}{%
\subsection{2. Theorem G1 and its
proof}\label{theorem-g1-and-its-proof}}

\emph{The alias sum.} With unitary Fourier transforms on the fine and
coarse tori, the coarse momenta \(K\in(-\pi,\pi]^D\) (units \(1/2a\))
receive the \(2^D\) fine momenta \(k_n=K/2+\pi n\), \(n\in\{0,1\}^D\),
and

\[(B\phi)^\wedge_{\mu\nu}(K)=2^{-D/2}\sum_nf_{\mu\nu}(k_n)\,\hat\phi_{\mu\nu}(k_n),\qquad
f_{\mu\nu}(k)=(1+e^{ik_\mu})(1+e^{ik_\nu}).\]

Hence the coarse covariance is
\(C'(K)_{\mu\nu,\rho\sigma}=2^{-D}\,t\sum_nf_{\mu\nu}(k_n)\overline{f_{\rho\sigma}(k_n)}\,P(k_n)_{\mu\nu,\rho\sigma}\).

\emph{The term \(n=0\).}
\(|f_{\mu\nu}(K/2)|^2=16\cos^2(K_\mu/4)\cos^2(K_\nu/4)=16(1+O(K^2))\),
and the lattice projector satisfies \(P(k)=P_{\rm cont}(\hat k)+O(k)\),
homogeneous of degree zero in the direction \(\hat k\); the \(O(k)\) is
a diagonal phase, and with plaquette-centred conventions (symbols
\(2\sin(k_j/2)\)) the error is \(O(k^2)\). So the \(n=0\) term is
\(2^{4-D}t\,[P_{\rm cont}(\hat K)+O(K)]\).

\emph{The terms \(n\ne0\).} If \(n_\mu=1\) or \(n_\nu=1\), then
\(1+e^{i(\pi+K_\mu/2)}=1-e^{iK_\mu/2}=O(K)\), so
\(f_{\mu\nu}(k_n)=O(K)\): a term whose two form factors are both
\(O(K)\) is \(O(K^2)\), and a mixed term with an \(O(1)\) partner is
covered by the closedness bound below. If \(n_\mu=n_\nu=0\) but
\(n_\lambda=1\) for some \(\lambda\notin\{\mu,\nu\}\), use closedness:
exact forms are closed, and the \((\mu\nu\lambda)\) component of
\(d\phi\) at \(k_n\) reads
\(d_\mu\hat\phi_{\nu\lambda}-d_\nu\hat\phi_{\mu\lambda}+d_\lambda\hat\phi_{\mu\nu}=0\)
with \(d_j=e^{ik_j}-1\), where \(d_\mu,d_\nu=O(K)\) and
\(d_\lambda=-2+O(K)\). So on the range of \(P(k_n)\) the component
\(\hat\phi_{\mu\nu}\) is \(O(K)\) times the others. Since \(P\) is an
orthogonal projector,
\(\|Pe_{\mu\nu}\|^2=(Pe_{\mu\nu})_{\mu\nu}=O(K)\|Pe_{\mu\nu}\|\), so
\(\|P(k_n)e_{\mu\nu}\|=O(K)\). An entry
\(P_{\mu\nu,\rho\sigma}=\langle Pe_{\mu\nu},Pe_{\rho\sigma}\rangle\) is
then \(O(K)\) times \(\|Pe_{\rho\sigma}\|\), which is itself \(O(K)\)
when \(n_\rho=n_\sigma=0\) (the same \(\lambda\) serves), while
otherwise \(f_{\rho\sigma}(k_n)=O(K)\) supplies the second factor. Every
term with \(n\ne0\) is \(O(K^2)\). \(\square\)

\emph{Iteration.} The proof used only that the input covariance is
supported on exact 2-forms (for closedness at \(k_n\)) and has the form
\(t[P_{\rm cont}+O(k)]\) at small \(k\) (for the \(n=0\) term). The
output has both properties, so by induction every step multiplies the
zero-momentum heat time by exactly \(2^{4-D}\). The four directional
Gaussian integrations of one step compose, in any order, to the
isotropic one (Fubini), so the exact directional defects of Proposition
2 and Corollary 2\(_s\) compose to a kernel with the same zero-momentum
coefficient; this is the free-field statement that their defects vanish
at zero momentum and constant flux.

\hypertarget{what-the-theorem-leaves-to-the-interacting-theory}{%
\subsection{3. What the theorem leaves to the interacting
theory}\label{what-the-theorem-leaves-to-the-interacting-theory}}

In the abelian theory the coupling has no loop corrections, and G1 is
the whole story at zero momentum. In the non-abelian theory the one-loop
matching coefficient \(c(\zeta,\mathcal K)\) of (15) depends on the full
kernel, and G1 localizes the obstruction: the finite endpoint term comes
from how the loop integrals feel the shape of the generated quadratic
kernel and from the generated vertices, never from a drift of the
tree-level coupling. Whether the shape parameters stay bounded, or drift
slowly, under the iteration is the question that the Remark (many steps)
of the four-dimensional note needs answered only in the weak form
\(o(n)\).

\hypertarget{consequence-for-state}{%
\subsection{4. Consequence for STATE}\label{consequence-for-state}}

Atlas cell 4: the Gaussian half of the composition obstruction is closed
at zero momentum (Theorem G1, iterated). The open half is the one-loop
dependence on the kernel shape and on generated vertices, with the
average rate requiring only \(c=o(n)\) along the iteration.
\end{document}
